% !TeX root = preliminaries.tex
\section{preliminaries}


\subsection{Varieties and rational maps}

Throughout this note, all schemes are of finite type and separated over a field or a noetherian integral scheme.
A variety is a geometrically integral scheme of finite type and separated over a field.
In other words, when base field is algebraically closed, a variety is an integral scheme of finite type and separated over the base field.
And in general case, $X$ is a variety if and only if its base change to the algebraic closure of the base field is a variety.
We usually use $\kk$ to denote an arbitrary field and $\kkk$ to denote an algebraically closed field.

When work over a non-algebraically closed field $\kk$, we often consider the set $X(\kkk)$ of $\kkk$-points of a variety $X$ over $\kk$ rather than the scheme-theoretic points of $X$.
Recall that for any field extension $\kk'/\kk$, the set of $\kk'$-points of $X$ is given by
\[ X(\kk') = \Mor_{\kk}(\Spec \kk', X) = \{ x: \Spec \kk' \to X \text{ over } \kk \}. \]
A point $x \in X(\kkk)$ is determined by the image of $x$ in $X$ and the induced embedding $\kappa(x) \hookrightarrow \kk'$ of the residue field $\kappa(x)$ of $x$ into $\kk'$.
For a scheme-theoretic closed point $\xi \in X$, we will say that $\xi \in X(\kk')$ if the residue field $\kappa(\xi)$ of $\xi$ is a subfield of $\kk'$, or equivalently, there exists $x \in X(\kk')$ such that the image of $x$ in $X$ is $\xi$.
Note that there may be at most $[\kappa(\xi):\kk]$ different points in $X(\kk')$ with the same image $\xi$ in $X$.
\begin{proposition}\label{prop:closed_points_on_varieties_over_non-algebraically_closed_fields}
	Let $X,Y$ be varieties over $\kk$ and $f: X \to Y$ be a morphism.
	Then
	\begin{enumerate}
		\item $X(\kkk)$ is canonically identified with the set of closed points of $X_\kkk$;
		\item $f: X(\kkk) \to Y(\kkk), x \mapsto f(x) = x \circ f$ coincides with the morphism $f_\kkk: X_\kkk \to Y_\kkk$ on the closed points under the above identification.
	\end{enumerate}
\end{proposition}
\begin{remark}
	If we consider the set of scheme-theoretic points of $X$ rather than $\kkk$-points, then the morphism $f: X \to Y$ may behave unlike our intuition.
	For example, let $X$ be a variety over finite field $\bbF_q$ and consider the absolute Frobenius morphism $\Frob_q$.
	Then $\Frob_q$ is the identity morphism on the set of scheme-theoretic points of $X$.
	But on $X(\overline{\bbF_q})$, it has many interesting dynamical properties.
\end{remark}

Let $X,Y$ be a variety over $\kk$, and $f: X \dashrightarrow Y$ be a rational map.
Recall that a rational map $f: X \dashrightarrow Y$ is a morphism defined on a dense open subset of $X$.
The fundamental method to study a rational map is considering its graph.
Let $\Gamma_f \subset X \times Y$ be the scheme-theoretic image of $(\id_U, f): U \to X \times Y$, where $U \subset X$ is the dense open subset where $f$ is defined.
This is equivalent to taking the closure of set-theoretic image of $(\id_U, f)$ in $X \times Y$.

\begin{definition}
	Let $X$ be a variety over $\kk$ and $f: X \dashrightarrow X$ be a dominant rational map.
	We say that a point $x \in X(\kkk)$ is a \emph{periodic point of $f$ (of period $r$)} if $f$ is defined at the orbit of $x$ and $f^r(x) = x$ for some $r > 0$.

	We say that a scheme-theoretic point $\xi \in X$ is a periodic point (of period $r$) of $f$ if it is the image of a periodic point (of period $r$) of $f$ in $X(\kkk)$.

	We denote by $\Per_r(f)$ the set of periodic points of $f$ of period $r$ and by $\Per(f) = \bigcup_{r > 0} \Per_r(f)$ the set of periodic points of $f$.

	A periodic point $x \in X(\kkk)$ of $f$ is called \emph{isolated} if it(or its image) is isolated in the $\Per_r(f)$ for its some period $r$.
\end{definition}

\begin{proposition}
	Let $X$ be a variety over $\kk$ and $f: X \dashrightarrow X$ be a dominant rational map.
	Then $\xi \in \Per_r(f)$ if and only if $f$ is defined at $\xi$ and $\xi$ lies in set-theoretic image of $\pi_1: \Gamma_{f^r} \cap \Delta_X \to X$ for some $r > 0$, where $\Gamma_{f^r}$ is the graph of $f^r$ and $\Delta_X$ is the diagonal of $X \times X$.
	% \Yang{}
\end{proposition}
\begin{proof}
	We have $\Per_r(f) \subset \pi_1(\Gamma_{f^r} \cap \Delta_X)$ since if $\xi \in \Per_r(f)$, then there exists $x \in X(\kkk)$ such that $\xi$ is the image of $x$ in $X$ and $f^r(x) = x$.
	Hence $(x,x) \in (\Gamma_{f^r} \cap \Delta_X)(\kkk)$ and $\xi$ is the image of $\pi_1(x,x) \in \pi_1(\Gamma_{f^r} \cap \Delta_X)(\kkk)$.

	Conversely, if $\xi \in \pi_1(\Gamma_{f^r} \cap \Delta_X)$, then there exists a scheme-theoretic point $\zeta \in \Gamma_{f^r} \cap \Delta_X$ such that $\pi_1(\zeta) = \xi$.
	Pick any $y \in \Gamma_{f^r} \cap \Delta_X(\kkk)$ with image $\zeta$ in $\Gamma_{f^r} \cap \Delta_X$, then $\pi_1(y)$ has image $\xi$ in $X$.
	Since $f$ is defined at $\xi$, we have $\pi_1(y)$ is a periodic point of $f$ of period $r$.
	Hence $\xi \in \Per_r(f)$.
	% \Yang{}
\end{proof}

The graph $\Gamma_f$ is a variety birational to $X$ via $\pi_1$ and $f$ can be recovered as the composition $\pi_2 \circ \pi_1^{-1}$, where $\pi_i: \Gamma_f \to X$ is the restriction of the $i$-th projection $X \times Y \to X$.
Using the graph, we can define the pullback and pushforward of divisors by $f$ as follows:
\[ f^*D := \pi_{1*} \pi_2^* D, \quad f_*D := \pi_{2*} \pi_1^* D. \]
However, not like morphisms, the pullback and pushforward by rational maps are not functorial in general, which causes some technical difficulties in the study of rational maps.
We will fix this issue by considering $\bfb$-cycles on birational models in \cref{subsec:intersection_product_on_birational_models}.

\begin{example}
	Consider the rational map $f: \bbP^2 \dashrightarrow \bbP^2$ defined by $[x:y:z] \mapsto [1/x:1/y:1/z] = [yz:xz:xy]$.
	Then graph of $f$ is given by
	\[ \Gamma_f = \{ ([x:y:z], [u:v:w]) \in \bbP^2 \times \bbP^2 : xu = yv = zw \}. \]
	Consider the first projection $\pi_1: \Gamma_f \to \bbP^2$ and $\bbA^2 \cong U \subset \bbP^2$ given by $z \neq 0$.
	Then
	\begin{align*}
		\pi_1^{-1}(U) & = \{ ([x:y:1], [u:v:w]) \in \Gamma_f : xu = yv = w \}   \\
		              & = \{ ([x:y:1], [u:v]) \in U \times \bbP^1 : xu = yv \},
	\end{align*}
	which is the blowup of $U$ at the origin.
	By similar argument, $\Gamma_f$ is the blowup of $\bbP^2$ at $[1:0:0]$, $[0:1:0]$ and $[0:0:1]$ and $\pi_1$ is the blowup morphism.

	We denote by $L_x$, $L_y$ and $L_z$ the coordinate lines in $\bbP^2$ given by $x=0$, $y=0$ and $z=0$ respectively,
	and $E_x$, $E_y$ and $E_z$ the exceptional divisors in $\Gamma_f$ over $[1:0:0]$, $[0:1:0]$ and $[0:0:1]$ respectively.
	On $\Gamma_f$, the strict transform of $L_x$'s and the exceptional divisors $E_x$'s form a hexagon of rational curves.
	Taking $H = L_x$ on $\bbP^2$, then we have $p_2^* H = \tilde{L_x} + E_y + E_z$.
	When we push forward by $\pi_1$, the strict transform $\tilde{L_x}$ is contracted to a point, while $E_y$ and $E_z$ are mapped isomorphically to $L_y$ and $L_z$ respectively.
	Hence we have $f^* H = \pi_{1*} \pi_2^* H = L_y + L_z$.
	Then we get $(f^*)^2 H \sim f^*2H \sim 4H$ while $(f^2)^*H = (\id)^* H = H$.
	% \Yang{}
	% On $\tilde{X} = \Gamma_f$, the induced rational map $\tilde{f}$ is a morphism given by $([x:y:z], [u:v:w]) \mapsto ([u:v:w], [x:y:z])$.
	% The Neron-Severi group $\NS(\tilde{X})$ is given by
	% \[ \bigoplus_{i=x,y,z} \bbZ \tilde{L_i} \oplus \bbZ E_i \big/ \langle \tilde{L_i} + E_j + E_k \sim \tilde{L_{i'}} + E_{j'} + E_{k'} | \{i,j,k\} = \{i',j',k'\} = \{x,y,z\} \rangle. \]
	% Then we have $\tilde{f}^* \tilde{L_x} = \tilde{L_x}$, $\tilde{f}^* E_x = E_x$, $\tilde{f}^* E_y = \tilde{L_y}$ and $\tilde{f}^* E_z = \tilde{L_z}$.
\end{example}

We need to consider family of dominant rational maps.
\begin{definition}
	Let $S$ be a noetherian integral scheme.
	A \emph{family of dominant rational maps} over $S$ is the following data:
	\begin{itemize}
		\item a scheme $X$ of finite type and flat over $S$ such that every fiber $X_s$ with $s \in S$ is a variety over $\kappa(s)$;
		\item an open dense subset $U \subseteq X$ and a morphism $f: U \to X$ over $S$ such that for every $s \in S$, the restriction of $f$ to the fiber $U_s$ gives a dominant rational map on $X_s$.
	\end{itemize}
	We denote by $f: X \dashrightarrow X$ the family of dominant rational maps over $S$.
\end{definition}

Given a family of dominant rational maps $f: X \dashrightarrow X$ over $S$, we can similarly define the graph $\Gamma_f \subset X \times_S X$ as the scheme-theoretic image (see \cite[Tag 01R5]{Stacks}) of the morphism $(i_U, f): U \to X \times_S X$, where $i_U: U \to X$ is the inclusion morphism.
As varieties, this is equivalent to taking the closure of set-theoretic image of $(\id_U, f)$ in $X \times_S X$.
We have $\Gamma_f$ is a scheme of finite type and separated over $S$ which is birational to $X$ via the first projection $\pi_1: \Gamma_f \to X$.

\begin{proposition}
	Let $f: X \dashrightarrow X$ be a family of dominant rational maps over $S$.
	Then for every $s \in S$, $\Gamma_{f_s}$ is a closed subscheme of $(\Gamma_f)_s$.
	%  \Yang{}
\end{proposition}
\begin{proof}
	The commutative diagram
	\[
		\begin{tikzcd}
			U_s \ar[r] \ar[d] & \Gamma_{f_s} \ar[r, hook] \ar[d,gray] & X_s \times_{\kappa(s)} X_s \ar[d] \\
			U \ar[r] & \Gamma_f \ar[r, hook] & X \times_S X
		\end{tikzcd}
	\]
	induces the gray arrow $\Gamma_{f_s} \to \Gamma_f$, and base changing by $s \to S$ gives the morphism $\Gamma_{f_s} \to (\Gamma_f)_s$.
	This morphism composing with the closed immersion $(\Gamma_f)_s \to X_s \times_{\kappa(s)} X_s$ yields the closed immersion $\Gamma_{f_s} \to X_s \times_{\kappa(s)} X_s$, so itself is a closed immersion.
\end{proof}

\begin{remark}
	In general, we have no equality $(\Gamma_f)_s = \Gamma_{f_s}$.
	For an example, see \cref{eg:non-continuity-of-mixed-degrees}
\end{remark}

The following lemma is well-known.
We will use it for lifting isolated periodic points from special fibers to the generic fiber.
\begin{lemma}\label{lem:intersection_dimension_estimate}
	Let $X$ be a noetherian scheme and $Y \subset X$ be a regular immersion.
	For any subscheme $Z \subset X$ and $x \in Y \cap Z$ (if $Y \cap Z \neq \emptyset$), we have
	\[ \dim_x (Y \cap Z) \geq \dim_x Y + \dim_x Z - \dim_x X. \]
\end{lemma}
\begin{proof}
	Suppose that $Y$ is locally defined by a regular sequence $f_1, \dots, f_r$ in the local ring $\calO_{X,x}$.
	Then we have
	\[ \dim_x (Y \cap Z) = \dim \calO_{Z,x} / (f_1, \dots, f_r) \geq \dim \calO_{Z,x} - r = \dim_x Z - (\dim_x X - \dim_x Y) \]
	by Krull's principal ideal theorem.
	% \Yang{}
\end{proof}

We may need the following special case of the flattening theorem of Gruson-Raynaud.
\begin{theorem}[{Gruson-Raynaud, ref. \cite[Tag 0815]{Stacks}}]\label{thm:Gruson-Raynaud_flattening}
	Let $S$ be a noetherian integral scheme and $X,Y$ of finite type and separated scheme over $S$.
	Let $f: X \to Y$ be a morphism over $S$ and $U \subseteq Y$ be an open subset such that $f^{-1}(U) \to U$ is flat.
	Then there exists a blowup $\pi: Y' \to Y$ with center in $Y \setminus U$ such that the induced morphism $f': X' \to Y'$ is flat, where $X'$ is the Zariski closure of $f^{-1}(U) \times_U \pi^{-1}(U)$ in $X \times_Y Y'$.
	%  \Yang{}
\end{theorem}




\subsection{Intersection product on birational models}\label{subsec:intersection_product_on_birational_models}

Let $X$ be a variety over $\kk$.
All birational models of $X$ form a directed system, we denote by $\frakX$.
Let $\CH^i(X)$ be the Chow group of cocycles of codimension $i$ on $X$; see \cite[Chapter 17]{Fulton13Intersection}.
% For our purpose, we only need the numerical equivalence classes of cocycles, namely $\upN^i(X)$; see \cite[Section 2]{Dang20degrees}.
Since the pullback by rational maps is not functorial, we need to consider the inductive limit of the Chow groups $\CH^i(X')$ for all birational models $X' \in \frakX$.
A dominant rational map $f: X \ratmap Y$ between varieties of the same dimension will induce a ``morphism'' $\frakX \to \frakY$ between the directed systems of birational models, which allows us to define a functorial pullback.

\begin{definition}
	We define
	\[ \CH^i(\frakX) := \varinjlim_{\pi:X'' \to X' \in \frakX} \left( \pi^* : \CH^i(X') \to \CH^i(X'') \right). \]
	An element in $\CH^i(\frakX)$ is called a \emph{$\bfb$-cocycle} of degree $i$ on $X$.
	It can be viewed as a pair $(X', \alpha' \in \CH^i(X'))$ modulo the equivalence relation generated by $(X_1, \alpha_1) \sim (X_2, \alpha_2)$ if there is a higher birational model $X'' \xrightarrow{\pi_i} X_i$ for $i=1,2$ such that $\pi_1^* \alpha_1 = \pi_2^* \alpha_2$.
	The model $X'$ is called a \emph{defining model} of $\alpha$.
	For $\alpha \in \CH^i(X)$, we still denote by $\alpha$ the $\bfb$-cocycle defined by $\alpha$.
	% We can define $\upN^i(\frakX)$ similarly by taking the inductive limit of $\upN^i(X')$.
	%  \Yang{}
\end{definition}

The intersection product $\CH^i(X') \otimes \CH^j(X') \to \CH^{i+j}(X')$ induces an intersection product
\[ \CH^i(\frakX) \otimes \CH^j(\frakX) \to \CH^{i+j}(\frakX) \]
by taking the intersection product on a common defining model.
It is well-defined since the pullback is a ring homomorphism.
The intersection number
\[ \CH^i(\frakX) \otimes \CH^{d-i}(\frakX) \to \bbZ \]
is also well-defined by the projection formula.

Let $f: X \dashrightarrow Y$ be a dominant rational map between varieties of the same dimension.
Then for every $\alpha \in \CH^i(\frakY)$, $f$ induces a rational map $X \ratmap Y'$ with $Y'$ a defining model of $\alpha$.
We will denote this data by $f: \frakX \to \frakY$.
Then $f$ induces a morphism $X' \to Y'$ for some birational model $X' \in \frakX$ by taking graph.
The pullback $f^* \alpha$ is a well-defined element in $\CH^i(X')$, and hence defines an element in $\CH^i(\frakX)$.
Easily check that the pullback $f^* \alpha$ is independent of the choice of the defining model $Y'$ of $\alpha$ and the choice of the birational model $X'$ of $X$.
In this way, if $f: X \dashrightarrow Y$ and $g: Y \dashrightarrow Z$ are two dominant rational maps between varieties of the same dimension, then we have $(g \circ f)^* = f^* \circ g^*$ as homomorphisms $\CH^i(\frakZ) \to \CH^i(\frakX)$.

Let us focus on the case of divisors, i.e. $i=1$.
In this case, note that the notion of nef, effective, big and pseudo-effective divisors are preserved by pullback, so we can also define these notions for $\bfb$-divisors on $\frakX$.
For ampleness, it is not preserved by pullback.
But we can still define an ample $\bfb$-divisor as a $\bfb$-divisor who has an ample defining model.

Let $f: X \dashrightarrow Y$ be a dominant rational map between varieties of the same dimension.
Then for every ample $\bfb$-divisor $L$ on $Y$, $f^* L$ will be also an ample $\bfb$-divisor on $X$.
Indeed, we may assume that $Y$ is a defining model of $L$ and $L$ is an ample divisor on $Y$.
Let $X'$ be the graph of $f$ and $\pi_1: X' \to X$ and $f': X' \to Y$ be the projections.
Then $f'$ is a generically finite morphism.
Taking the Stein factorization of $f'$, we get a factorization $X' \xrightarrow{\pi'} X'' \xrightarrow{f''} Y$ with $\pi'$ birational and $f''$ finite.
Then $f''^* L$ is an ample divisor on $X''$ since $f''$ is finite.
This gives an ample defining model $X''$ of $f^* L$.

Let $Z \subset X$ be a closed subvariety of dimension $i$ and $\alpha \in \CH^i(\frakX)$.
We can take a defining model $X'$ of $\alpha$ which is higher than $X$ by $\pi: X' \to X$.
Suppose that $Z$ is not contained in the indeterminacy locus of $\pi^{-1}$, then we can define the intersection number
\[ \alpha \cdot Z := \alpha_{X'} \cdot \tilde{Z}, \]
where $\tilde{Z}$ is the strict transform of $Z$ on $X'$.
Under suitable assumptions making the intersection numbers defined, we still have the projection formula for pullback by rational maps.

\begin{proposition}\label{prop:projection_formula_for_pullback_by_rational_maps}
	Let $f: X \ratmap Y$ be a dominant rational map between spaces of the same dimension, $\alpha \in \CH^i(Y)$ and $Z \subset X$ be a closed subvariety of dimension $i$.
	Suppose that $Z$ is not contained in the indeterminacy locus of $f$, then the intersection number $f^* \alpha \cdot Z$ is defined and we have
	\[ f^* \alpha \cdot Z = \alpha \cdot f_*Z, \]
	where $f_*Z := [\fkk(Z) : \fkk(\overline{f(Z)})] \cdot \overline{f(Z)}$.
	%  \Yang{}
\end{proposition}
\begin{proof}
	Let $X'$ be the graph of $f$ and $\pi_1: X' \to X$ and $f': X' \to Y$ be the projections.
	Note that the indeterminacy locus of $f$ is the same as the indeterminacy locus of $\pi_1^{-1}$, so $Z$ is not contained in the indeterminacy locus of $\pi_1^{-1}$.
	By definition, we have $X'$ is a defining model of $f^* \alpha$ since $f': X' \to Y$ is a morphism.
	Let $\tilde{Z}$ be the strict transform of $Z$ on $X'$, then we have
	\[ f^* \alpha \cdot Z = \alpha_{X'} \cdot \tilde{Z} = f'^* \alpha \cdot \tilde{Z} = \alpha \cdot f'_* \tilde{Z} = \alpha \cdot f_*Z. \]
\end{proof}




\subsection{Siu's inequality and dynamical degrees}

First we list some forms of Siu's inequality which will be used in the proof of the main theorem also in the proof of the properties of dynamical degrees.

Then original version of Siu's inequality is the following.
\begin{theorem}[{Siu's criterion, cf. \cite[Theorem 2.2.15]{Laz04a}}]\label{thm:Siu_inequality}
	Let $X$ be a projective variety of dimension $d$ and $L,M$ be ample divisors on $X$.
	Then
	\[ M - \varepsilon \frac{1}{d} \frac{M^d}{L \cdot M^{d-1}} L \]
	is big for every $\varepsilon \in (0,1)$.
	%  \Yang{}
\end{theorem}

For higher codimensional cocycles, we have the following generalization of Siu's inequality.
Note that this is a little weaker than the original version of Siu's inequality since bigness is replaced by the weaker condition that the intersection number is positive.

\begin{theorem}[{Generalized Siu's inequality, ref. \cite{JiangLi2023Siu-inequality}}]\label{thm:generalized_Siu_inequality}
	Let $X$ be a projective variety of dimension $d$ and $\alpha \in \CH^k(X)$, $\beta \in \CH^{d-k}(X)$ be intersections of nef divisors on $X$ and $H$ a nef divisor on $X$.
	Then we have
	\[ (\alpha \cdot \beta) \cdot H^{d} \leq \binom{d}{k} (\alpha \cdot H^{d-k}) (\beta \cdot H^{k}). \]
\end{theorem}

By restricting the above inequality to hyperplane sections and taking limit to nef divisors, we can get the following version.

\begin{corollary}
	Let $X$ be a projective variety of dimension $d$ and $\alpha \in \CH^k(X)$, $\beta \in \CH^l(X)$ and
	$\eta \in \CH^{d-k-l}(X)$ be intersections of nef divisors on $X$ and $H$ be a nef divisor on $X$.
	Then we have
	\[ (\alpha \cdot \beta \cdot \eta) \cdot (H^{k+l} \cdot \eta) \leq \binom{k+l}{k} (\alpha \cdot H^{l} \cdot \eta) (\beta \cdot H^{k} \cdot \eta). \]
\end{corollary}

Now we can show the existence of dynamical degrees and their log-concavity.

\begin{proposition}
	Let $f: X \dashrightarrow X$ be a rational map on a projective variety $X$ of dimension $d$.
	Then the limit
	\[ \lambda_k(f) := \lim_{n \to \infty} \left(\deg_{k,H}(f^n)\right)^{1/n} \]
	exists and is independent of the choice of the ample divisor $H$.
\end{proposition}
\begin{proof}
	First

	\Yang{}
\end{proof}

\begin{proposition}
	The sequence of dynamical degrees $\lambda_k(f)$ is log-concave, i.e.
	\[ \lambda_k(f)^2 \geq \lambda_{k-1}(f) \lambda_{k+1}(f)\]
	for every $1 \leq k \leq d-1$.
	Or equivalently, the sequence of Lyapunov exponents $\mu_k(f)$ is non-increasing, i.e.
	\[ \mu_k(f) \geq \mu_{k+1}(f) \]
	for every $1 \leq k \leq d-1$.
\end{proposition}
\begin{proof}

	\Yang{}
\end{proof}




\subsection{Semi-continuous functions on noetherian schemes}

Let $S$ be an integral noetherian scheme, and $\theta: S \to \bbR$ be a function.
We say that $f$ is \emph{lower semi-continuous} if for every $r \in \bbR$, the set $\{ s \in S : f(s) > r \}$ is open in $S$.

\begin{remark}
	The limit of a sequence of lower semi-continuous functions may not be lower semi-continuous.
	\Yang{}
\end{remark}

\begin{lemma}\label{lem:characterization_of_lower_semi-continuous_functions_on_noetherian_schemes}
	A function $\theta: S \to \bbR$ is lower semi-continuous if and only if the following condition holds:
	\begin{enumerate}
		\item for every $\eta \in S$ and every specialization $\xi \in \overline{\{\eta\}}$, we have $\theta(\xi) \leq \theta(\eta)$;
		\item for every $\eta \in S$ and $a < \theta(\eta)$, there is an open subset $U \subseteq \overline{\{\eta\}}$ such that $\theta(\xi) > a$ for every $\xi \in U$.
	\end{enumerate}
\end{lemma}
\begin{proof}
	\Yang{}
\end{proof}


\begin{lemma}\label{lem:properties_of_lower_semi-continuous_functions_on_noetherian_schemes}
	Let $\theta: S \to \bbR$ be a lower semi-continuous function.
	Then
	\begin{enumerate}
		\item $\theta$ is continuous with respect to the constructible topology on $S$;
		\item if $\theta$ is discrete, then its image is finite;
		\item let $\eta \in S$ be the generic point, then $\lim_{s \to \eta} \theta(s) = \theta(\eta)$.
	\end{enumerate}
\end{lemma}
\begin{proof}
	\Yang{}
\end{proof}




\subsection{The Frobenius}

In this subsection, we review some fundamental properties of the Frobenius morphism, which may be trivial for experts.
But for me, they are not so trivial at least at the beginning of this note.

\begin{definition}
	Let $S$ be a scheme over $\bbF_p$.
	Then the \emph{absolute Frobenius} morphism $\Frob_{S,p}: S \to S$ is the morphism defined locally by the $p$-th power map on the structure sheaf $\calO_S$.
	%  \Yang{}
\end{definition}

\begin{proposition}
	Let $S,T$ be schemes over $\bbF_p$ and $f: S \to T$ be a morphism.
	Then $f$ is commutative with the absolute Frobenius $\Frob_p$.
\end{proposition}

\begin{definition}
	Let $S$ be a scheme over $\bbF_p$ and $X$ be a scheme over $S$.
	We denote by $X^{(p)}$ the base change of $X$ along the absolute Frobenius $\Frob_{S,p}$, i.e. $X^{(p/S)} = X \times_{S, \Frob_{S,p}} S$, called the \emph{Frobenius twist} of $X$.
	It is equipped with an $S$-scheme structure given by the second projection $X^{(p/S)} \to S$.

	The absolute Frobenius $\Frob_{S,p}: S \to S$ induces an $S$-morphism $\id_X \times \Frob_{S,p}: X \to X^{(p/S)}$, called the \emph{relative Frobenius} morphism of $X$ over $S$ and denoted by $\Frob_{X/S, p}$.
	In a diagram, we have
	\[ \begin{tikzcd}
			X \arrow[rrd, bend left, "\Frob_{X,p}"] \arrow[rdd, bend right] \arrow[rd, "\Frob_{X/S,p}"] & & \\
			& X^{(p/S)} \arrow[r] \arrow[d] & X \arrow[d] \\
			& S \arrow[r, "\Frob_{S,p}"] & S.
		\end{tikzcd} \]
	%  \Yang{}
\end{definition}
\begin{remark}
	We often omit the subscript $S$ on the Frobenius twist when there is no confusion, i.e. $X^{(p)}$ for $X^{(p/S)}$.
	We also often omit the subscript $X$ on the Frobenius morphism when there is no confusion, i.e. $\Frob_p$ for $\Frob_{X,p}$ and $\Frob_{/S,p}$ for $\Frob_{X/S,p}$.
	And we write $\Frob_{p^n}$ (resp. $\Frob_{/S,p^n}$) for the $n$-th iterate of $\Frob_p$ (resp. $\Frob_{/S,p}$).
\end{remark}

\begin{example}
	Let $S = \Spec \kkk$ be the spectrum of an algebraically closed field of characteristic $p$ and $X \subset \bbA^n_\kkk$ be an affine variety defined by $f_1, \dots, f_m \in \kkk[t_1, \dots, t_n]$.
	Then the Frobenius twist $X^{(p)}$ is given by $X^{(p)} = \Spec \kkk[t_1, \dots, t_n]/(f_1^{(p)}, \dots, f_m^{(p)})$, where $f_i^{(p)}$ is obtained by applying the $p$-th power map to the coefficients of $f_i$.
	The projection $X^{(p)} \to X$ is given by $(a_1, \dots, a_n) \mapsto (a_1^{1/p}, \dots, a_n^{1/p})$, while the relative Frobenius $\Frob_{/S,p}: X \to X^{(p)}$ is given by $(a_1, \dots, a_n) \mapsto (a_1^p, \dots, a_n^p)$.
	% \Yang{}
\end{example}

\paragraph{Relative Frobenius as endomorphisms when defined over finite fields}
The most interesting case for us is when $S = \Spec \bbF_q$ or $S = \Spec \overline{\bbF_q}$ and $X$ is a variety over $S$.
First consider the case when $S = \Spec \bbF_q$.
Then the absolute Frobenius $\Frob_q: \Spec \bbF_q \to \Spec \bbF_q$ is the identity morphism, so the Frobenius twist $X^{(q)}$ is just $X$ and the relative Frobenius $\Frob_{/S,q}$ is just the absolute Frobenius $\Frob_q$.

More generally, let $\kkk$ be an arbitrary algebraically closed field of characteristic $p$ and $X$ be a variety over $\kkk$.
Suppose that $X$ is defined over some finite field $\bbF_q$.
That is, there exists a variety $X_0$ over $\bbF_q$ and an isomorphism $X \cong X_0 \times_{\bbF_q} \kkk$.
Then we have diagrams
\[
	\begin{tikzcd}
		X \arrow[r] \arrow[rdd] \arrow[rd, red] & \Spec \kkk \arrow[rd, "\Frob_q"] \arrow[rdd] & \\
		& X \arrow[r] \arrow[d] & \Spec \kkk \arrow[d] \\
		& X_0 \arrow[r] & \Spec \bbF_q
	\end{tikzcd}
	\quad \leadsto \quad
	\begin{tikzcd}
		X \arrow[rrd, red] \arrow[rdd] \arrow[rd, blue] & & \\
		& X^{(q)} \arrow[r] \arrow[d] & X \arrow[d] \\
		& \Spec \kkk \arrow[r, "\Frob_q"] & \Spec \kkk.
	\end{tikzcd}
\]
In particular, this gives an isomorphism $X^{(q)} \cong X$ (the blue arrow) over $\kkk$.
In general, for a variety $X$ over $\kkk$, we still have an isomorphism $X^{(q)} \cong X$ as schemes but not as $\kkk$-schemes.
Using this $\kkk$-scheme isomorphism, the relative Frobenius $\Frob_{/\kkk,q}: X \to X^{(q)}$ can be viewed as a morphism $X \to X$, which is different from the absolute Frobenius $\Frob_q: X \to X$.
Another more direct way to get the relative Frobenius is to take the product
\[ \Frob_{/\kkk,q} = (\Frob_{q})_\kkk = \Frob_{X_0/\bbF_q, q} \times \id_\kkk: X_0 \times_{\bbF_q} \kkk \to X_0 \times_{\bbF_q} \kkk. \]
This is the construction we will use in the following.

If $\kkk = \overline{\bbF_q}$, then every variety over $\kkk$ is defined over some finite field, so the above construction applies to every variety over $\kkk$.

\begin{example}[Graph of Frobenius on $\bbA^1_{\bbF_p}$]
	Let $X = \bbA^1_{\bbF_p} = \Spec \bbF_p[t]$ be the affine line over $\bbF_p$.
	We consider the absolute Frobenius $F = \Frob_p: X \to X$ given by $\bbF_p[t] \to \bbF_p[t], f \mapsto f^p$.
	Note that even though $F$ is the identity on the underlying topological space of $X$, the graph $\Gamma_F$ of $F$ is not the diagonal in $X \times X$.

	Since $(\id,F): X \to X \times X$ is given by $\bbF_p[t] \otimes_{\bbF_p} \bbF_p[t] \to \bbF_p[s], f(t) \otimes g(s) \mapsto f(t) g(t^p)$, the graph $\Gamma_F$ is given by $\Spec \bbF_p[s,t]/(s - t^p)$.
	This is different from the diagonal $\Delta_X$ given by $\Spec \bbF_p[s,t]/(s-t)$.

	The reason is that the underlying space $|X \times_{\bbF_p} X|$ of the fiber product $X \times_{\bbF_p} X$ is not the same as the product $|X| \times |X|$ of the underlying spaces.
	For example, let $\alpha \in \overline{\bbF_p} \setminus \bbF_p$ with minimal polynomial $f(t) \in \bbF_p[t]$.
	This defines a closed point $\frakm_\alpha = (f(t)) \in |X|$.
	But there are several closed points in $|X \times_{\bbF_p} X|$ mapping to $\frakm_\alpha$ via both projections.
	For example, I claim that the ideal $\frakm_n = (f(t), s-t^{p^n})$ defines a closed point in $|X \times_{\bbF_p} X|$ mapping to $\frakm_\alpha$ via both projections for every $n \geq 0$.
	Before proving the claim, note that $\frakm_0 \neq \frakm_1$.
	Otherwise, we would have $t-t^p \in \frakm_0 \cap \bbF_p[t] = (f(t))$, which is impossible since $f(t)$ is irreducible of degree at least $2$ and $t-t^p$ factors as $t(t-1)(t-2) \cdots (t-(p-1))$.

	We denote by $\phi_1: \bbF_p[t] \to \bbF_p[t,s]$ and $\phi_2: \bbF_p[s] \to \bbF_p[t,s]$ the homomorphisms corresponding to the two projections $X \times_{\bbF_p} X \to X$.
	Then we have $\pi_i(\frakm_n) = \phi_i^{-1}(\frakm_n) = \ker(\bbF_p[t] \xrightarrow{\phi_i} \bbF_p[t,s]/\frakm_n)$.
	Factor $\bbF_p[t,s] \to \bbF_p[t,s]/\frakm_n$ as $\bbF_p[t,s] \to \bbF_p[t] \to \bbF_p[t]/(f(t))$, where the first map is given by $s \mapsto t^{p^n}$ and the second map is the quotient map.
	Then $\pi_1(\frakm_n) = (f(t)) = \frakm_\alpha$.
	For every $g(s) \in \bbF_p[s]$, regard $g(s)$ as a polynomial in $\bbF_p(t)[s]$, we have $s-t^{p^n} \mid g(s) - g(t^{p^n})$ in $\bbF_p(t)[s]$.
	Then there exists $h(t,s) \in \bbF_p[t,s]$ such that $g(s) - g(t^{p^n}) = h(t,s)(s-t^{p^n})$.
	In particular, $g(s) \in \frakm_n$ if and only if $g(t^{p^n}) = g(t)^{p^n} \in \frakm_n$, and the latter is equivalent to $g(t) \in (f(t)) = \frakm_\alpha$ since $\frakm_\alpha$ is a prime ideal.
	Hence $\pi_2(\frakm_n) = (f(s)) = \frakm_\alpha$ as well.
\end{example}

\begin{proposition}
	Let $X$ be a variety over $\bbF_{q}$.
	Then for every closed point $x \in X(\overline{\bbF_q})$, $x \in X(\bbF_{q^n})$ if and only if $\Frob_{q}^n(x) = x$.
	In particular, every closed point of $X$ is a periodic point of the Frobenius $\Frob_q$.
\end{proposition}
\begin{proof}
	Let $x$ be given by a morphism $x: \Spec \kk \to X$ with $\kk$ the residue field of $x$.
	Note that $\Frob_q^n(x)$ is given by the diagram
	\[ \begin{tikzcd}[column sep=large]
			\Spec \kk \arrow[r, "x"] \arrow[rd, "\Frob_q^n(x)" description] \arrow[d,"\Frob_q^n"'] & X \arrow[d, "\Frob_q^n"] \\
			\Spec \kk \arrow[r, "x"] & X.
		\end{tikzcd} \]
	Hence $\Frob_q^n(x) = x$ if and only if $\Frob_q^n: \Spec \kk \to \Spec \kk$ is the identity morphism, which is equivalent to $\kk \subset \bbF_{q^n}$.
\end{proof}

\begin{proposition}
	Let $f: X \to Y$ be a morphism of varieties over $\kk$.
	Suppose that $f,X,Y$ are defined over some finite field $\bbF_q$.
	Then $f$ is commutative with the relative Frobenius $\Frob_{/\kk,q}$.
	%  \Yang{}
\end{proposition}
\begin{proof}
	Choose $f_0: X_0 \to Y_0$ defined over $\bbF_q$ such that $f = f_0 \times_{\bbF_q} \kk$.
	Since $f_0$ is commutative with the absolute Frobenius $\Frob_q$, the base change $f = f_0 \times_{\bbF_q} \kk$ is also commutative with the relative Frobenius $\Frob_{/\kk,q} = (\Frob_q)_\kk$.
\end{proof}
